\chapter{{\fontsize{14pt}{21pt}\selectfont\MakeUppercase{Тестирование, развертывание и внедрение системы}}}
\label{cha:testing}

\section{Тестирование системы}

Тестирование выполнено для подтверждения соответствия реализованного решения функциональным и нефункциональным требованиям, сформулированным в главе~1. Проверка включала серверную часть, клиентский интерфейс и интеграционные пользовательские сценарии.

Тестирование серверной части

Для backend-части (Django + DRF) использован модульный подход \cite{django_docs_2026,drf_docs_2026}:
\begin{compactlist}
  \item проверка корректности доменных моделей и ограничений целостности;
  \item проверка прав доступа по ролевой модели \texttt{user/moderator/admin};
  \item проверка сценариев модерации, обработки жалоб и уведомлений;
  \item проверка валидации входных данных и обработки ошибок API.
\end{compactlist}

Критериями приемки серверной части являются:
\begin{compactlist}
  \item успешное применение миграций и отсутствие конфликтов схемы БД;
  \item корректная работа основных API-методов (CRUD и служебные actions);
  \item отсутствие критических ошибок в системной проверке Django.
\end{compactlist}

Тестирование клиентской части

Для frontend-части (React + Effector) проведено функциональное тестирование ключевых пользовательских потоков \cite{react_docs_2026,effector_docs_2026}:
\begin{compactlist}
  \item авторизация и работа с профилем пользователя;
  \item просмотр и фильтрация каталогов рецептов и планов;
  \item формирование плана питания в конструкторе;
  \item работа с отзывами, избранным и списками покупок;
  \item обработка ошибок запросов и отображение уведомлений.
\end{compactlist}

Отдельно проверены граничные условия конструктора плана питания:
\begin{compactlist}
  \item длительность плана задается в диапазоне от 1 до 30 дней;
  \item шаг изменения длительности составляет 1 день;
  \item количество перекусов задается в диапазоне от 0 до 5 в сутки.
\end{compactlist}

Интеграционное тестирование

Интеграционная проверка выполнена на уровне бизнес-сценариев:
\begin{compactlist}
  \item регистрация пользователя и заполнение профиля;
  \item создание рецепта, отправка на модерацию, изменение статуса;
  \item формирование плана питания и генерация списка покупок;
  \item публикация отзыва, создание жалобы, получение уведомления;
  \item обращение к ИИ-ассистенту через серверный прокси.
\end{compactlist}

Результаты тестирования подтвердили работоспособность основных сценариев использования системы и согласованность работы клиентской и серверной частей.

\section{Развертывание на сервере}

Сборка и контейнеризация серверной части

Для унификации окружения и воспроизводимости развертывания приложение упаковано в Docker-контейнеры \cite{docker_docs_2026}. Проект включает:
\begin{compactlist}
  \item контейнер базы данных PostgreSQL;
  \item контейнер backend-сервиса (Django + Gunicorn);
  \item контейнер frontend-сервиса (Vite + React).
\end{compactlist}

Конфигурация backend-контейнера приведена в листинге~\ref{listing:backend-dockerfile}.

\begin{listing}[H]
\caption{Dockerfile серверной части}
\label{listing:backend-dockerfile}
\begin{mdframed}
\begin{minted}[breaklines=true]{bash}
FROM python:3.13-slim
ENV PYTHONDONTWRITEBYTECODE=1 PYTHONUNBUFFERED=1
WORKDIR /app
COPY requirements.txt /app/requirements.txt
RUN pip install --no-cache-dir -r /app/requirements.txt
COPY . /app
EXPOSE 8000
CMD ["gunicorn", "config.wsgi:application", "--bind", "0.0.0.0:8000"]
\end{minted}
\end{mdframed}
\end{listing}

Конфигурация frontend-контейнера приведена в листинге~\ref{listing:frontend-dockerfile}. Сборка выполняется через \texttt{npm ci}, а запуск dev-сервера --- через Vite \cite{npm_docs_2026,vite_docs_2026,vite_deploy_docs_2026}.

\begin{listing}[H]
\caption{Dockerfile клиентской части}
\label{listing:frontend-dockerfile}
\begin{mdframed}
\begin{minted}[breaklines=true]{bash}
FROM node:22-alpine
WORKDIR /app
COPY package.json package-lock.json /app/
RUN npm ci
COPY . /app
EXPOSE 5173
CMD ["npm", "run", "dev", "--", "--host", "0.0.0.0", "--port", "5173"]
\end{minted}
\end{mdframed}
\end{listing}

Оркестрация сервисов

Для запуска полного стенда используется \texttt{docker-compose.yml}, включающий базы данных, backend и frontend (листинг~\ref{listing:compose-file}) \cite{docker_compose_docs_2026}.

\begin{listing}[H]
\caption{Фрагмент docker-compose.yml}
\label{listing:compose-file}
\begin{mdframed}
\begin{minted}[breaklines=true]{yaml}
services:
  db:
    image: postgres:16-alpine
  backend:
    build: ./Backend
    ports: ["8000:8000"]
    command: >
      sh -c "python manage.py migrate &&
             gunicorn config.wsgi:application --bind 0.0.0.0:8000"
  frontend:
    build: ./Frontend
    ports: ["5173:5173"]
\end{minted}
\end{mdframed}
\end{listing}

Порядок запуска

Последовательность развертывания:
\begin{compactlist}
  \item перейти в корневую директорию проекта;
  \item скачать (при необходимости обновить) образы командой \texttt{docker compose pull};
  \item выполнить сборку контейнеров командой \texttt{docker compose build};
  \item запустить окружение командой \texttt{docker compose up -d};
  \item проверить доступность сервисов: frontend (\texttt{http://localhost:5173}), backend API (\texttt{http://localhost:8000/api/health/}).
\end{compactlist}

Такой подход сокращает время подготовки среды и обеспечивает идентичное поведение приложения на разных хостах.

\section{Перспективы развития системы}

Дальнейшее развитие FoodPlanner целесообразно выполнять по следующим направлениям:
\begin{compactlist}
  \item расширение доменной персонализации (сезонность, бюджет, расширенные медицинские ограничения);
  \item развитие рекомендательных алгоритмов и повышение качества ИИ-консультаций;
  \item интеграция с внешними сервисами (трекеры активности, носимые устройства, продуктовые каталоги);
  \item расширение аналитики пользовательского прогресса и поведенческих метрик;
  \item усиление контуров качества: автоматизированные e2e-тесты, CI/CD и централизованная телеметрия.
\end{compactlist}

\section{Выводы по главе}

В третьей главе представлены результаты тестирования, подтверждающие работоспособность ключевых пользовательских и служебных сценариев системы. Показано развертывание приложения в контейнерной среде Docker с использованием единого файла оркестрации сервисов. Также определены приоритетные направления дальнейшего развития, направленные на повышение качества персонализации, масштабируемости и эксплуатационной зрелости платформы.